% !TeX root = ../main.tex
\chapter{Glossary}
\label{app:glossary}

The entries are terms whose meaning is specific to underground work and is
not recoverable from general engineering usage. Symbols and units are in the
List of Symbols, and typical parameter values in
Appendix~\ref{app:constants}.

\begin{description}
\item[Advance core] The volume of ground immediately ahead of the face, whose stiffness and strength govern face stability and pre-convergence.
\item[Convergence] Reduction in the cross-sectional dimensions of an excavation following stress redistribution.
\item[Disturbance factor ($\Dfac$)] The Hoek--Brown parameter representing excavation-induced damage and stress relaxation adjacent to the opening, from 0 for undisturbed mechanical excavation to 1 for heavy blast damage. It applies to the damaged annulus only and not to the whole rock mass (Chapter~\ref{ch:rock-mass-parameters}).
\item[Equivalent dimension ($D_e$)] The excavation span or height divided by the excavation support ratio, used to enter the $Q$-system support chart.
\item[Excavation support ratio (ESR)] A factor, between about 0.5 and 5, by which the span or height of an excavation is divided before the $Q$-system support chart is entered. It expresses the degree of security required according to use and design life, and it runs inversely to it: temporary openings take the highest values and the most important permanent caverns the lowest, so that a more important excavation presents a larger equivalent dimension and attracts heavier support (Table~\ref{tab:esr}).
\item[Extrusion] Longitudinal deformation of the advance core along the tunnel axis, through the face and into the excavation, as distinct from the radial pre-convergence of the profile ahead of it (Chapter~\ref{ch:construction}).
\item[Hogging and sagging] The two senses of curvature imposed on a structure by a settlement trough: hogging, convex upward, over the outer part of the trough, and sagging, concave upward, over its centre (Chapter~\ref{ch:movement-monitoring}).
\item[Ground reaction curve (GRC)] The relationship between a notional internal pressure, reduced from the in-situ value towards zero, and the resulting radial convergence. It describes the response of the ground alone and is independent of the support actually installed (Chapter~\ref{ch:ccm}).
\item[Longitudinal deformation profile (LDP)] The relationship between convergence and distance from the face, used to determine how much deformation has already occurred when support is installed.
\item[Overbreak] Excavation beyond the design profile, whether caused by blast damage, structurally controlled fallout or machine overcut.
\item[Pre-convergence] Convergence of the theoretical profile ahead of the face.
\item[Squeezing] Time-dependent large deformation of weak ground under stress exceeding its strength.
\item[Stand-up time] The period for which an unsupported excavation span remains stable.
\item[Reinforced ribs of sprayed concrete (RRS)] Arches of sprayed concrete reinforced with bars, used in the heavier support categories of the $Q$-system (Chapter~\ref{ch:empirical-design}).
\item[Support characteristic curve (SCC)] The relationship between support pressure and convergence for an installed support system.
\item[Volume loss ($V_L$)] The volume of ground movement at the surface expressed as a proportion of the excavated volume.
\end{description}

\chapterreferences
